\part{XXZ / Vertex Models}

\section{Quantum groups}

\begin{exercise}[The open XXZ chain and Hopf algebra symmetry]
Let $V=\C^{\{+,-\}}$ be the physical arrow space of the section
Row-to-row transfer matrices and let $q>0$, $q\ne1$,
$d=q+q^{-1}$. On $V\otimes V$ define a kernel $e$ that vanishes
on equal arrows and has mixed-arrow kernel
\[
 \begin{pmatrix}q&-1\\-1&q^{-1}\end{pmatrix}
\]
in the order $(+-,-+)$. On $V^{\otimes N}$ put
$H^{\rm open}=-J\sum_{i=1}^{N-1}e_i/2$, $J>0$.
An algebra action on a tensor product is specified by a coproduct:
a bialgebra is a unital associative algebra $A$ with algebra maps
$\Delta:A\to A\otimes A$ and $\varepsilon:A\to\C$ satisfying
coassociativity and
$(\varepsilon\otimes\id)\Delta=\id=(\id\otimes\varepsilon)\Delta$.
A Hopf algebra also has an antipode $S$ satisfying
$m(S\otimes\id)\Delta=m(\id\otimes S)\Delta=\eta\varepsilon$.
\begin{enumerate}[label=(\alph*)]
\item Show that $e$ is positive, $e^2=de$, and
\[
 H^{\rm open}=J\sum_{i=1}^{N-1}
 \left(S_i^xS_{i+1}^x+S_i^yS_{i+1}^y
 +\frac d2(S_i^zS_{i+1}^z-\tfrac14)\right)
 -\frac{J(q-q^{-1})}{4}(S_1^z-S_N^z).
\]
Use the spin conventions of the section Yang--Baxter equation.
\item For a Hopf algebra prove that
$x\cdot(v\otimes w)=\Delta(x)(v\otimes w)$ defines an action
and that iterating $\Delta$ is independent of parentheses.
Prove uniqueness of $S$ as the inverse of $\id$ under convolution
$f*g=m(f\otimes g)\Delta$.
\item On a finite-dimensional module define
$(xf)(v)=f(S(x)v)$ on its dual. Prove that evaluation and the
coevaluation of the section Transfer matrices are module maps,
using the antipode and snake identities.
\end{enumerate}
\end{exercise}

\begin{exercise}[Monoidal categories and quasitriangular symmetry]
A category consists of objects, morphism sets, associative
composition, and identity morphisms. A monoidal category has a
tensor bifunctor, a unit, an associator
$\alpha_{U,V,W}:(U\otimes V)\otimes W\to U\otimes(V\otimes W)$,
and left and right unit isomorphisms. Their coherence equations
are the pentagon
\[
 \alpha_{U,V,W\otimes X}\alpha_{U\otimes V,W,X}
 =(\id_U\otimes\alpha_{V,W,X})\alpha_{U,V\otimes W,X}
                  (\alpha_{U,V,W}\otimes\id_X)
\]
and the triangle equating the two unit-removal maps on
$(U\otimes\mathbf1)\otimes V$.
A braiding consists of natural isomorphisms
$c_{U,V}:U\otimes V\to V\otimes U$ satisfying, with associators
inserted, the two hexagon identities
\[
 c_{U\otimes V,W}=(c_{U,W}\otimes\id_V)(\id_U\otimes c_{V,W}),
 \quad
 c_{U,V\otimes W}=(\id_V\otimes c_{U,W})(c_{U,V}\otimes\id_W).
\]
\begin{enumerate}[label=(\alph*)]
\item Verify these structures for vector spaces and for modules
over a bialgebra, using its coproduct for the latter.
\item A quasitriangular Hopf algebra has an invertible
$\mathcal R\in A\otimes A$ with
$\mathcal R\Delta(x)=\Delta^{\rm op}(x)\mathcal R$,
$(\Delta\otimes\id)\mathcal R=\mathcal R_{13}\mathcal R_{23}$,
and $(\id\otimes\Delta)\mathcal R=\mathcal R_{13}\mathcal R_{12}$.
Prove that $c_{U,V}=\tau\mathcal R$ gives a braiding on its
modules and derive the Yang--Baxter identity.
\item Prove the braid relation for $c$ directly from its
hexagons and naturality. Check that a cocommutative Hopf algebra
has $\mathcal R=1$ and the usual tensor symmetry.
\end{enumerate}
\end{exercise}

\section{\texorpdfstring{$U_q(\mathfrak{sl}_2)$}{Uq(sl2)}}

\begin{exercise}[The quantum enveloping algebra and its fundamental spin action]
For $q\in\C^\times$, $q^2\ne1$, define $U_q(\mathfrak{sl}_2)$
by generators $E,F,K,K^{-1}$ and relations
\[
 KEK^{-1}=q^2E,\quad KFK^{-1}=q^{-2}F,\quad
 [E,F]=\frac{K-K^{-1}}{q-q^{-1}},\quad KK^{-1}=K^{-1}K=1.
\]
Set
\[
 \Delta(E)=E\otimes K+1\otimes E,\quad
 \Delta(F)=F\otimes1+K^{-1}\otimes F,\quad
 \Delta(K)=K\otimes K,
\]
$\varepsilon(E)=\varepsilon(F)=0$, $\varepsilon(K)=1$,
$S(E)=-EK^{-1}$, $S(F)=-KF$, $S(K)=K^{-1}$.
Write $[n]_q=(q^n-q^{-n})/(q-q^{-1})$.
\begin{enumerate}[label=(\alph*)]
\item Verify that these maps respect all the defining relations
and give a Hopf algebra. Prove $S^2(x)=KxK^{-1}$.
\item On the physical arrow space let $K$ multiply the two
arrow states by $q,q^{-1}$, let $E$ flip $-$ to $+$ and
annihilate $+$, and let $F$ flip $+$ to $-$ and annihilate $-$.
Verify the relations and prove that $e$ commutes with the
coproduct action on two sites.
\item By induction using the commutator relation prove
\[
 [E,F^r]=[r]_q F^{r-1}
 \frac{q^{1-r}K-q^{r-1}K^{-1}}{q-q^{-1}}.
\]
Define a type-$1$ module to be a module with a direct
decomposition into $K$-eigenspaces of weights in $q^\Z$.
Show that $E$ and $F$ move these spaces by weights $+2$ and $-2$.
\end{enumerate}
\end{exercise}

\begin{exercise}[Quantum symmetric powers and highest-weight modules]
Assume $q$ is not a root of unity. In the physical tensor space
$V^{\otimes n}$ define $V_n=\bigcap_{i=1}^{n-1}\ker e_i$ and
$V_0=\C$.
\begin{enumerate}[label=(\alph*)]
\item In the sector with $r$ down arrows show that the equations
$e_i\psi=0$ relate values at configurations differing by an
adjacent exchange by a factor of $q$ or $q^{-1}$.
Check consistency by counting inversions of the down-arrow set.
Deduce that this sector of $V_n$ is a line for $0\le r\le n$.
\item Prove that $V_n$ is a submodule. If $L$ is its
highest-weight line, prove that $F^r:L\to V_n[q^{n-2r}]$
is an isomorphism for $0\le r\le n$, and that
$EF^r|_L=[r]_q[n-r+1]_qF^{r-1}|_L$.
Deduce irreducibility and $\dim V_n=n+1$.
\item Let a finite-dimensional simple type-$1$ module have
highest nonzero weight space $L$. Apply the same commutator
identity at the first vanishing lowering map to prove that its
highest weight is $q^n$ for some $n\ge0$ and that the module
is isomorphic to $V_n$.
\item Compute the intrinsic weighted trace
$\operatorname{tr}_q(f)=\Tr(Kf)$ on $V_n$ and obtain
$\operatorname{tr}_q(\id_{V_n})=[n+1]_q$.
Prove cyclicity for module maps and multiplicativity under
tensor products.
\end{enumerate}
\end{exercise}

\section{Hecke / Temperley--Lieb algebras}

\begin{exercise}[Braid, Hecke, and Temperley--Lieb relations]
The braid group $B_N$ has generators $\sigma_i$ with
$\sigma_i\sigma_{i+1}\sigma_i=\sigma_{i+1}\sigma_i\sigma_{i+1}$
and commuting generators for $|i-j|\ge2$.
The Hecke algebra $H_N(q)$ imposes additionally
$(g_i-q)(g_i+q^{-1})=0$.
The Temperley--Lieb algebra $TL_N(d)$ has generators $e_i$ with
\[
 e_i^2=de_i,\qquad e_ie_{i\pm1}e_i=e_i,\qquad
 [e_i,e_j]=0\quad(|i-j|\ge2).
\]
\begin{enumerate}[label=(\alph*)]
\item Verify these relations for the two-arrow kernel $e$
of the section Quantum groups, by checking the three-site
arrow configurations.
\item In $TL_N(q+q^{-1})$ put $g_i=q-e_i$.
Prove the braid and Hecke relations and derive
$g_i^{-1}=g_i-(q-q^{-1})$. Deduce the homomorphism
$H_N(q)\to TL_N(d)$.
\item Define $TL_N(d)$ also by planar noncrossing pairings
of $N$ input and $N$ output endpoints, with multiplication
by stacking and with each closed loop replaced by $d$.
Reduce an arbitrary pairing using an innermost arc and
prove that the generators satisfy the presentation above:
straightening adjacent cups and caps gives the three relations,
and the relations reduce every word to a scalar times a word
\[
 (e_{i_1}e_{i_1-1}\cdots e_{j_1})\cdots
 (e_{i_r}e_{i_r-1}\cdots e_{j_r}),
 \quad i_1<\cdots<i_r,\quad j_1<\cdots<j_r,\quad j_a\le i_a.
\]
Here $1\le j_a\le i_a<N$, and the empty word is allowed.
Match these words bijectively with the planar pairings.
Count pairings by the arc at the first boundary point to derive
the Catalan recurrence and $\dim TL_N(d)=\binom{2N}{N}/(N+1)$.
\end{enumerate}
\end{exercise}

\begin{exercise}[Baxterization of the Temperley--Lieb interaction]
Let $q=e^\eta$ and $d=2\cosh\eta$. In $TL_N(d)$ put
$X_i(u)=\sinh(\eta-u)\id+\sinh u\,e_i$.
\begin{enumerate}[label=(\alph*)]
\item Expand both sides and use the Temperley--Lieb relations
to prove
\[
 X_i(u)X_{i+1}(u+v)X_i(v)
 =X_{i+1}(v)X_i(u+v)X_{i+1}(u).
\]
Reduce the remaining scalar condition to hyperbolic addition.
\item Prove $X_i(0)=\sinh\eta\,\id$ and
$X_i(u)X_i(-u)=(\sinh^2\eta-\sinh^2u)\id$.
Calculate the logarithmic derivative at zero and identify
its nonconstant part with $e_i/\sinh\eta$.
\end{enumerate}
\end{exercise}

\section{Representation theory}

\begin{exercise}[Intrinsic isotypic decomposition and Clebsch--Gordan fusion]
Let $q$ be not a root of unity. The Casimir element is
$C=FE+(qK+q^{-1}K^{-1})/(q-q^{-1})^2$.
A module is semisimple if it is a direct sum of simple modules.
\begin{enumerate}[label=(\alph*)]
\item Check that $C$ is central and acts on $V_n$ by
$c_n=(q^{n+1}+q^{-n-1})/(q-q^{-1})^2$.
Show that the $c_n$ are distinct for $n\ge0$.
\item For a finite-dimensional type-$1$ module decompose by
generalized Casimir eigenspaces. In a summand whose simple
quotients are $V_n$, show that its weight-$q^n$ space generates
a sum of copies of $V_n$: the lowering maps are injective
until level $n$, and level $n+1$ is absent.
Compare dimensions with a composition series to prove
complete reducibility.
\item Prove that the evaluation map gives the canonical
decomposition
\[
 W\simeq\bigoplus_{n\ge0}V_n\otimes\Hom_{U_q}(V_n,W).
\]
Compute multiplicities by subtracting consecutive weight
dimensions, and derive
$V_m\otimes V_n\simeq\bigoplus_{r=0}^{\min(m,n)}V_{m+n-2r}$.
\item Apply this to $V^{\otimes N}$ and obtain the multiplicity
$\binom Nr-\binom N{r-1}$ of $V_{N-2r}$.
Show that the open-chain Hamiltonian acts only on the
multiplicity spaces, so that every energy in such a space
has multiplicity at least $N-2r+1$.
\end{enumerate}
\end{exercise}

\begin{exercise}[The classical spin limit]
Let $\mathfrak{sl}_2$ be the trace-zero endomorphisms of a
two-dimensional vector space with bracket $[A,B]=AB-BA$.
Its enveloping algebra is the tensor algebra modulo
$x\otimes y-y\otimes x-[x,y]$; it acts on tensor powers by
the coproduct $x\mapsto x\otimes1+1\otimes x$.
\begin{enumerate}[label=(\alph*)]
\item Set $q=e^\eta$ and write $K=e^{\eta H}$ on the
weight spaces of $V_n$. Take $\eta\to0$ in the ladder
relations and obtain $[H,E]=2E$, $[H,F]=-2F$, $[E,F]=H$.
Identify $V_n$ with $\Sym^n(V)$ by its tensor action.
\item Put $S^z=H/2$, $S^x=(E+F)/2$,
$S^y=(E-F)/(2\ii)$, and $s=n/2$.
Compute the central operator
$(S^x)^2+(S^y)^2+(S^z)^2=s(s+1)\id$ and the character
$\Tr e^{tS^z}=\sinh((2s+1)t/2)/\sinh(t/2)$.
Use the induced inner product on $\Sym^n(V)$ to verify
self-adjointness of the spin operators.
\end{enumerate}
\end{exercise}

\section{Rational/trigonometric hierarchy}

\begin{exercise}[Trigonometric transfer operators and the rational limit]
Use the trigonometric $R(u)$ and $H_{\Delta,h}$ of the
sections Spectral parameter and Yang--Baxter equation.
\begin{enumerate}[label=(\alph*)]
\item Under $u=\epsilon x$, $\eta=\epsilon$ prove
$\epsilon^{-1}R(u)\to x\id+P$ and $\Delta\to1$.
Pass to the limit in the Yang--Baxter and RTT relations.
\item Show that the classical limit of the open-chain kernel
is $e=\id-P$. Derive
$H^{\rm open}=\frac J2\sum_i(P_{i,i+1}-\id)$ and its
global rotation symmetry.
\item On a periodic chain derive the one-magnon energy
$h+J(\cos k-1)$ at $\Delta=1$ and compare with the
trigonometric band $h+J(\cos k-\cosh\eta)$.
Keep the different open and periodic boundary terms explicit.
\end{enumerate}
\end{exercise}

\begin{exercise}[Prediction: quantum-group multiplets of an open XXZ segment]
For $q>0$, $q\ne1$, take the Hermitian open Hamiltonian
$H^{\rm open}=-J\sum_i e_i/2$ with the boundary field fixed
in the section Quantum groups.
\begin{enumerate}[label=(\alph*)]
\item For two sites use $e^2=de$ and the intrinsic decomposition
$V\otimes V=V_2\oplus V_0$ to derive a triplet at energy zero
and a singlet at energy $-Jd/2=-J\Delta$.
\item For $N$ sites prove that any simple energy in the
multiplicity space of $V_{N-2r}$ has exactly $N-2r+1$
states unless its energy coincides with that of another
multiplet. Show that it has one state in each magnetization
sector $S^z_{\rm tot}=(N-2r)/2-m$ for $0\le m\le N-2r$.
\item For an operator in the $V_2$ component of
$\End(\mathcal H_N)$, with action
$x\cdot O=\sum x_{(1)}OS(x_{(2)})$, apply Clebsch--Gordan
fusion to derive the selection rule
$V_n\to V_{n-2}\oplus V_n\oplus V_{n+2}$, omitting
negative labels. Prove that a scalar sector $V_0$ couples
only to $V_2$ through such an operator.
\item State the two-site absorption energy $J\Delta$, the
multiplet degeneracies, and the selection rule as separate
spectroscopic predictions. Identify the boundary field,
anisotropy, and transformation law of the probe required
to test them.
\end{enumerate}
\end{exercise}
